\section{Diseño experimental, semillas y condiciones}
\label{sec:diseno}

\subsection{Diseño factorial}
Se cruzan dos horarios (base y candidata congelada) con tres niveles de demanda:
0.80, 1.00 y 1.20 veces el volumen de referencia. El diseño tiene \NConditions{} condiciones
y cada una usa \NReps{} réplicas independientes, para \NRuns{} ejecuciones finales. Los
niveles 0.80 y 1.20 funcionan también como análisis de sensibilidad alrededor de la
condición nominal.

\subsection{Semillas y congelamiento de la candidata}
\begin{itemize}
  \item \textbf{Optimización:} semillas \OptSeedFirst--\OptSeedLast{} (diez), solo para
  buscar y seleccionar la candidata; semilla del algoritmo \AlgSeed.
  \item \textbf{Evaluación:} semillas \EvalSeedFirst--\EvalSeedLast{} (\NReps{}), no usadas
  para buscar, seleccionar ni modificar la candidata.
  \item \textbf{Emparejamiento:} dentro de cada nivel de demanda, la réplica $i$ de la base y
  la de la candidata comparten semilla (números aleatorios comunes); réplicas distintas usan
  semillas distintas.
\end{itemize}
El archivo \texttt{frozen\_candidate.json} se escribe antes de evaluar las semillas finales.

\subsection{Estadísticos}
Por condición y métrica se reportan el número de réplicas, la media, la desviación estándar
muestral, el intervalo de confianza normal del 95\,\% para la media, la mediana y los
percentiles 10 y 90 entre réplicas. Para comparar horarios se usa la diferencia pareada
$\text{base}-\text{candidata}$ (positiva significa menos espera con la candidata), su media,
su intervalo normal del 95\,\% y el número de parejas favorables. El análisis es
\emph{descriptivo}: no se ejecuta una prueba de hipótesis y no se usa la palabra
``significativo''.

\subsection{Reproducción}
Todo el directorio \texttt{evidence/final/} se regenera con
\texttt{python scripts/generate\_evidence.py}; la evaluación sola, con la candidata ya
congelada, con \texttt{python scripts/run\_experiments.py --candidate
evidence/final/frozen\_candidate.json}. El manifiesto registra versiones
(Python \PythonVersion), sumas de verificación de los datos procesados, configuración,
semillas y comandos. La generación usada en este informe corresponde al commit
\texttt{\GitCommit}.
